\section{Model}
The following section develops a formal model to analyze interactions between political elites deciding whom to endorse in a primary election. Given that endorsements influence candidate viability, elites must weigh their ideological preferences against the political benefits of backing a winning candidate. The model introduces uncertainty in the nomination process and allows for variation in the structure of political rewards, which will dictate whether endorsements act as strategic complements or substitutes.

\subsection*{Endorsements \& Electoral Outcomes}
Consider two political elites that must decide whether to endorse the establishment candidate or an outsider as their party's nominee in a general election. Each candidate $c \in \{E, O\}$ has an exogenous and known quality $v_c \in [0,1]$, measuring their underlying electability. While candidate quality reflects characteristics such as competence, charisma, and overall appeal to the general electorate, elite endorsements shape intra-party competition by signaling viability and influencing voter expectations.

Endorsements contribute to candidate viability, but their impact is subject to diminishing marginal returns. The probability that candidate $E$ wins the nomination is determined by a function $f(n)$, where $n$ is the number of endorsements received, as well as some exogenous uncertainty. 

Rather than modeling explicitly the complexity of voter updating in a multi-endorser setting, the autonomic nomination rule encapsulates the reality that endorsements systematically increase a candidate's likelihood of being nominated. This approach simplifies the analysis while still capturing the core strategic incentives facing elites. Modeling explicit voter learning introduces additional complexity without altering the fundamental insight that endorsements influence nominations in a probabilistic manner. Instead, the nomination rule serves as a reduced-form representation of the effect of endorsements, shifting focus away from the intricate mechanics of voter responses.

Let the establishment candidate win the party nomination if
\[f(n) + V = \frac{\alpha n}{1+ \beta n} + V > 0\]

where $V\sim \mathcal{N}(v_E - v_O, \sigma^2)$, with cumulative distribution $F_V(\cdot)$. Here, $\alpha >0$ represents the baseline effectiveness of endorsements in shaping candidate viability, while $\beta > 0$ governs the diminishing rate of returns. This functional form ensures that the marginal benefit of each additional endorsement decreases as more endorsements accrue. Early endorsements play a crucial role in shaping voter perceptions and media narratives, while later endorsements add relatively less informational value. By incorporating diminishing returns, the model captures a more realistic scenario where early fragmentation may occur before eventual coordination. This formulation prevents implausible outcomes where a candidate’s probability of nomination approaches one simply due to an arbitrarily large number of endorsements.

The uncertainty generated by $V$ accounts for real-world factors such as polling inaccuracies, voter preference shifts, and unforeseen campaign events. Even with strong elite backing, primary elections are inherently volatile—debate performances, media coverage, fundraising surges, and scandals can all impact nomination outcomes. If outcomes were fully deterministic, elites would have no incentive to strategically weigh their endorsements. By incorporating uncertainty, the model reflects the reality that endorsements are made in a probabilistic environment where risk-averse elites may hesitate to commit to a candidate who, although appealing ideologically, may be perceived as electorally weak.

Once a nominee is selected, their chances of winning the general election depend on their individual quality and external factors beyond elite control. Similar to the nomination process, the general election is modeled as a probabilistic event. Conditional on being the party nominee, candidate $c$ wins the general election if and only if

\[v_c + \epsilon > 0\]

where $\epsilon \sim \mathcal{N}(0, \sigma_{\epsilon}^2)$. Thus, the probability that $c \in \{E, O\}$ wins the general election when they are the party nominee is

\[Pr(v_c + \epsilon > 0) = 1-Pr(v_c + \epsilon \leq 0) = 1 -\Phi\bigg(-\frac{v_c}{\sigma_{\epsilon}}\bigg) = \Phi\bigg(\frac{v_c}{\sigma_{\epsilon}}\bigg) \equiv p_c\]

This formulation reflects the uncertainty of general elections even after the primary, as factors like economic conditions, national political trends, and opposition candidate quality can influence voters. Unlike the nomination process, endorsements do not directly influence the general election. Instead, the nominee’s pre-existing candidate quality determines their electability. This distinction underscores why elites must carefully weigh both ideological alignment and electability concerns when deciding whether to endorse a candidate.

\subsection*{Political Rewards \& Strategic Incentives}
Elites' endorsement decisions depend on the political rewards they expect to gain from backing a particular candidate. These rewards embody tangible political benefits such as positions in the future administration (e.g., cabinet appointments or ambassadorships), access to party resources, greater policy influence, or support from party leadership in subsequent elections. To capture how reward structures impact endorsement incentives, the model considers two distinct settings:

\begin{enumerate}
\item \textbf{Fixed Rewards:} Every elite who endorses the winning candidate receives the same reward, independent of how many elites endorse the same candidate. This structure is best exemplified in the tit-for-tat nature of endorsements among elected officials. Newly elected presidents often wield their influence and commit resources to their supporters during subsequent midterm elections or gubernatorial campaigns. Under fixed rewards, elites have a strong incentive to coordinate their support, as aligning with the frontrunner maximizes the potential of receiving such political benefits.

\item \textbf{Split Rewards:} In this setting, rewards are split evenly among all endorsers of the winning candidate. Thus, the presence of multiple endorsers reduces the value each elite receives from supporting the eventual nominee. Examples include limited prestigious appointments such as cabinet posts that become diluted or less valuable when many elites compete for them. In this environment, elites strategically differentiate their endorsements to maximize individual influence or secure greater personal returns by backing less popular, yet still viable candidates.
\end{enumerate}

Formally, elite utility depends on both ideological alignment with candidates and expected political returns. Let $e_i \in \{E, O\}$ denote the candidate elite $i \in \{1,2\}$ endorses. Ideological alignment with candidates is represented by $x_i \in [-1,1]$, measuring elite $i$'s benefit from endorsing the establishment candidate. Elite $i$'s utility is given by

\[u_i(e_1, e_2; v_E, v_O) = x_i\cdot \mathbbm{1}\{e_i = E\} + Pr(e_i \text{ nominated}|n)\cdot \Phi(\frac{v_{e_i}}{\sigma_{\epsilon}})\cdot R\]

The first term captures the ideological payoff of endorsing $E$. The second term represents the expected political rewards, factoring in how endorsements impact a candidate's likelihood of nomination. Under the fixed reward structure, $R$ remains constant for each endorser. Under the split reward structure, total rewards are divided equally among endorsers: if only one elite endorses the winning candidate, they receive the full reward ($R=1$); if both endorse the winner, each receives half ($\frac{1}{2}R$).

Elites announce their endorsements simultaneously after observing candidate qualities and electoral uncertainty. Their endorsements reflect ideological preferences as well as strategic calculations about endorsement competition or coordination, depending critically on the nature of political rewards available.
